\documentclass[12pt]{article}
\usepackage[utf8]{inputenc}
\usepackage[T1]{fontenc}
\usepackage{hyperref}
\title{Saho Language Resources for K-12 Curriculum Development}
\author{Compiled by Claude Sonnet 5.5 (claude-sonnet-5-5) for Andr\'as Kornai}
\date{}
\begin{document}
\maketitle

\section{Introduction}
Saho (ISO 639-3: ssy) is an Afro-Asiatic (Cushitic) language of Eritrea and Ethiopia. Banti and Vergari estimate about 190,000 native speakers in Eritrea, citing Lewis (2009). Saho is one of Eritrea's nine instructional languages.

\section{Primary Educational Resources}
\begin{itemize}
\item Global Partnership for Education blog, ``Eritrea: educating children in their mother tongue''. It describes free textbooks written in Saho for Hutsit Primary School, supported by GPE and UNICEF. \url{https://globalpartnership.org/blog/eritrea-educating-children-their-mother-tongue}
\end{itemize}

\section{Linguistic and Reference Materials}
\begin{itemize}
\item Banti and Vergari, ``The Saho of Eritrea and the Documentation of their Language and Cultural Heritage'' (PDF). It mentions the Eritrean Ministry of Education's Saho primary schools in Gash Barka. \url{https://www.sahoarchive.org/wp-content/uploads/2014/08/Banti_Vergari_The_Saho_of_Eritrea_and_the_Documentation_of_their_Language_and_Cultural_Heritage-libre.pdf}
\item ``A sketch of Saho grammar'' (Banti and Vergari). It notes the 2003 Vergari and Vergari Basic Saho-English-Italian Dictionary, published in Asmara and used in Saho elementary schools. \url{https://www.academia.edu/758286/A_sketch_of_Saho_grammare}
\item Saho Archive, grammatical notes. \url{https://www.sahoarchive.org/grammatical-notes/}
\end{itemize}

\section{Religious and Community Resources}
\begin{itemize}
\item Saho Archive, a documentation project on the Saho language and cultural heritage. \url{https://www.sahoarchive.org/?p=9}
\end{itemize}

\section{Digital and Audio Resources}
The Saho Archive site above may hold audio and text material. This was not verified in the search results.

\section{Curriculum Development Context}
Saho is used for instruction in Eritrea's early grades, and primary textbooks and a school dictionary exist, according to the sources above. Eritrean Ministry of Education materials should be sought directly. The search did not locate them online.

\end{document}